\def\velkostpisma{9}
\def\velkostriadku{12}
\input makra.tex % nacitanie Ivanom pripravenych nastaveni a prikazov
\hyphenation{star-šov-stvo} % rozdelenie slov na konci riadku, treba tu uviest slova, ktore sam nepozna

\spravodaj{3}{2026}


\clanok{{\it ČO NÁS ČAKÁ V~MARCI}}


\clanok{Verš na mesiac}
V~marci sa budeme učiť nový verš, ktorý dostala pre rok~2026 malá besiedka:

„Či neviete, že ste chrám Boží a Duch Boží prebýva vo vás?“ (1K~3,16)


\clanok{Výročné zborové členské zhromaždenie}
Výročné zborové členské zhromaždenie sa bude konať v~nedeľu 8.~3.~2026 o~15.30~hod. v~modlitebni na Palisádach. Súčasťou bohatého programu sú aj voľby ďalšieho kazateľa pre náš zbor. Členovia zboru boli už podrobne informovaní.


\clanok{Rozhovory na Palisádach}
V~sobotu, dňa 14.~3.~2026 o~14.00~hod. na Palisádach prebehne štvrtý z~„Rozhovorov na Palisádach“. Hosťom bude predseda ICEJ Slovensko Peter Švec s~témou „Prečo je dobré rozumieť Izraelu?“ Pozvánka s~viac informáciami bude čoskoro aj vo vašich e-mailových schránkach. Ak vás táto téma zaujíma, ste srdečne vítaní. Pozvite aj svojich priateľov.


\clanok{Sestry}
\begitems
* Všetky sestry pozývame na stretnutie Svetového dňa modlitieb, ktorý sa bude konať prvý marcový piatok -- 6.~3.~2026 o~17.00~hod., v~Cirkvi bratskej na Cukrovej~4. Spolu s~ostatnými kresťanmi na svete sa budeme modliť za Nigériu.
* Stretnutie sestier nášho zboru sa uskutoční v~stredu 25.~3. o~17.30~hod. na Zrínskeho~2. Tému a hosťa oznámime čoskoro.
* 1.~--~3.~5.~2026 sa bude konať v~Olomouci Konferencia pre ženy BJB ČR a SR. Konferencia nesie názov „Milovaná“. Registrovať sa môžete do 29.~3.~2026 na webovej stránke \ulink[https://www.konferenceprozeny.cz/registrace/]{konferenceprozeny.cz/registrace}.
\enditems


\clanok{Pôstne zamyslenia}
Tento mesiac sme vo veľkonočnom pôstnom období. Vo štvrtky (12.~3., 19.~3., 26.~3. a 2.~4.) v~čase biblických hodín o~18.00~hod. na Palisádach máme pripravené kontemplatívne bohoslužby -- Pôstne zamyslenia. Sprevádzať nás nimi budú bratia Filip~Barkóczi a Dávid Mário Chuchút. Všetci ste srdečne vítaní.


\clanok{Veľkonočný koncert}
Veľký spevokol a komorný orchester nášho zboru nás pozýva na Koncert Veľkonočných piesní. Prvý je v~sobotu 21.~3.~2026 o~17.00~hod. v~našej modlitebni na Palisádach a druhý v~nedeľu 22.~3.~2026 o~17.00~hod. v~ECAV kostole Sv. Trojice v~Petržalke na Strečnianskej~15.


\clanok{{\it OHLIADNUTIE SA ZA FEBRUÁROM}}


\clanok{Správy zo staršovstva}
Vo februári sa staršovstvo stretlo dvakrát.

Na prvom stretnutí nás s.~R.~Krajčiová informovala o~jej misijnej ceste do Filipín a spolu sme sa modlili za jej ochranu a aby všetci účastníci tejto misijnej akcie boli Jeho nástrojmi pri ohlasovaní evanjelia. So s.~Ľubkou Kohútovou sme prešli výsledné hospodárenie zboru za rok~2025 a upravený návrh zborového rozpočtu na rok~2026. Okrem malých zmien týkajúcich sa členenia rozpočtu je navrhnutý nový technický fond, ktorý bude slúžiť na prevádzku a obnovu techniky, ktorú využívame počas bohoslužieb. Sme vďační Pánovi, že nás drží a požehnáva aj v~oblasti financií a že sme ako zbor mohli v~minulom roku podporiť aj cez mimoriadne zbierky viac akcií v~sume takmer 24~000~€. Sestra Ľubka nám opäť pripomenula, že bude potrebná jej náhrada v~tejto službe
(v~ktorej je verne už mnoho rokov). Aj touto cestou vás bratia a sestry, ktorí máte k~tomuto (účtovníctvo) potrebné skúsenosti, oslovujeme, aby ste rozmýšľali o~vašej účasti v~tejto dôležitej službe. Ďalším hosťom na staršovstve bol brat Janko Štefko. Informoval nás o~potrebe obmeny v~radoch diakonov a tiež o~krokoch, ktoré pre to urobil. Sme radi, že viacerí z~vás už prejavili záujem byť v~tejto službe.

Ďalším bodom bola otázka prenájmu priestorov v~prízemí na Zrínskeho (bývalá kancelária zboru) E.~Kakašovej ako ateliér. Z~dôvodu vyťaženosti hospodárskeho výboru z~projektov, ktoré budú prebiehať tento rok (obnova fasády na Palisádach, čistička odpadových vôd na Chvojnici) a nedostatku financií na obnovu priestorov starej kancelárie, sme súhlasili s~ročným prenájmom do konca februára~2027. Potom by sme radi tieto priestory zrekonštruovali, aby mohli slúžiť zboru.

Brat R.~Nemec predstavil ideu spoločných obedov v~neďalekých priestoroch Strediska evanjelickej diakonie. A~venovali sme sa ešte našim návrhom kandidátov do Rady BJB. K~tomuto účelu bolo v~januári aj jedno mimoriadne stretnutie staršovstiev západnej oblasti na Súľovskej.

Druhé stretnutie staršovstva malo viacerých hostí. Prišla s.~Eva Molnárová, ktorá prejavila záujem o~členstvo v~našom zbore a povedala nám svoje svedectvo, ako ju Pán Ježiš zachránil. Sestra N.~Hovorková a br. Ľ.~Kešjar využili príležitosť otvoreného staršovstva vo veci návrhu nového rozpočtu zboru. Sme vďační za ich pripomienky. A~tiež bol medzi nami br. P.~Žembery, aby sme pripravili veci týkajúce sa voľby ďalšieho kazateľa zboru, ktoré sa začnú na výročnom zborovom zhromaždení 8.~3.~2026. Prosím, zapojme sa do modlitieb, aby sme v~tejto voľbe boli Božím nástrojom v~presadení Jeho vôle.
\autor{Ľubomír Syč}
\vfill\break


\clanok{Lyžovačka}
Aj tento rok sme niektorí z~nás počas jarných prázdnin vycestovali na sever \smiley. Rodiny, jednotlivci, snáď všetky vekové kategórie sme si nenechali ujsť našu už tradičnú zborovú „lyžovačku“.
Prežili sme opäť pekný oddychový týždeň obklopený horskou prírodou v~Račkovej doline. Mnohí sme sa poriadne vylyžovali, vysánkovali, či poprechádzali.
Oddýchli sme si pri spoločných rozhovoroch, hrách či pripravených večerných zamysleniach a ranných modlitbách. Sme vďační Pánu Bohu za každoročnú možnosť byť takto spolu mimo mesta.
\autor{Katarína Kerekréty}


\clanok{Očami Slávky}
Raz mi moja ómama povedala: „Slávinka, ako by som včera chodila do školy...“ Vtedy som tomu nerozumela, ale dnes už áno.
S~láskavým úsmevom sa často pozerám do minulosti a vybavujem si spomienky spred mnohých rokov. A~pred pár dňam moje spomienky ožili ako v~rozprávkovej Narnii.
Cesta vlakom s~vnúčatami do Liptovského Mikuláša, autobusom do Pribyliny a už len kúsok do milovanej Račkovej. Bola som tam posledne pred mnohými rokmi, keď boli moje deti o~čosi staršie ako dnes vnúčatká. Hanka a Dušan sú stále rovnakí milí a láskaví, chata má nový odev a prekrásnu jedáleň s~nádhernými preskleným výhľadom do celej doliny.
Každé ráno stretnutia pri šálke kávy, chlebíku vo vajci, sladkých koláčikoch... spoločné ranné úvahy a stíšenie, potom rozpŕchnutie na svahy a poobede znova chata. Krik a smiech detí, spoločné hry, rozhovory, klavír s~ľuďmi, ktorých sme dlho nevideli, s~ktorými často len v~nedeľu prehodíme pár slov. Večerné hry, rozhovory, smiech....nové priateľstvá a milé prekvapenia, že nové kamarátky sú vlastne sesternice z~tretieho kolena. Drahý Nebeský Otče, ďakujem, ani neviem nájsť slová vďačnosti za to, aký si ku nám dobrý.
\autor{Slávka Kráľová}
\vfill\break


\clanok{Mládežnícka konferencia}
„Príď zažiť, čo naozaj znamená byť živí. Semináre. Športy. Hry. Worship. Ľudia ako ty.“ To boli heslá v~pozvánke na tohtoročnú Konferenciu mládeže BJB v~Banskej Bystrici.

V~čase uzávierky spravodaja sme stihli položiť niekoľkým mládežníkom iba zopár otázok. Ak máte záujem sa dozvedieť viac, pýtajte sa osobne, určite vám/nám radi porozprávajú.

\begitems \style n
* Je niečo konkrétne, čo si očakával/očakávala od mládežníckej konferencie?
\begitems
* A.~D.: Všetko, čo som očakávala, tak aj bolo. Veľmi som sa tešila na túto ml. konferenciu.
* B.~K.: Veľa zábavy, smiechu, únavy a krásne prežitý čas s~Bohom a priateľmi.
* R.~N.: Srandu.
\enditems
* Na ktorom si bol seminári si bol/bola čo ti z~neho najviac utkvelo v~pamäti?
\begitems
* A.~D.: Bola som na seminári u~Dávida Chuchúta. Jeho téma bola, prečo zostať na Slovensku a prečo odísť a bola som tam, lebo chcem odísť.
* B.~K.: „Zostať alebo odísť? Slovensko ako krajina pre moju budúcnosť“ od D.~M.~Chuchúta. Išla som, pretože ma zaujímalo, kam by som mala ísť na výšku a či sa oplatí zostať v~malej krajine -- Slovensko.
* R.~N.: Bol som na Dávidovom (Chuchútovom) o~tom, či zostať na Slovensku alebo nie a to preto, lebo by som chcel v~buducnosti odísť študovať alebo žiť do zahraničia.
\enditems
* Čo všetko si si z~konferencie odniesol/odniesla?
\begitems
* A.~D.: Očakávam, že získam inšpiráciu a povzbudenie, aby som viac prežívala Božiu prítomnosť v~každodennom živote a zakúsila Jeho lásku v~plnej miere.
* B.~K.: Super, určite pôjdem aj na budúci rok, veľmi som si to užila.
* R.~N.: Fúha, no bolo to super. Páčilo sa mi tam dobrá nálada a zobral som si z~toho asi, že na Slovensku nie je tak zle, a že aj západ má svoje chyby a treba sa za to modliť. Ale to s~tým modlením neviem, či povedal, ale ja som si to tak zobral.
\enditems
\enditems
\vfill\break

\clanok{Zbierky v~januári}
\table{lr}{
Na misiu			& 547~€ \cr
Na investičný fond	& 417~€ \cr
}
\vskip1em

Pripomíname, že zbierky sa u~nás konajú pravidelne takto:
\vskip-1ex\begitems
* každú 2. nedeľu v~mesiaci je zbierka venovaná misii a
* každú 4. nedeľu je zbierka na investície.
\enditems

{\bf Pre dosporenie na rekonštrukciu} sme vypočítali príspevok na člena vo výške 150~€. Prispieť je možné formou zbierky počas 4.~nedieľ v~mesiaci alebo bankovým prevodom na zborový účet s~uvedením variabilného symbolu 777, prípadne aj osobne s.~Ľubke Kohútovej.

Bankové spojenie: SK36 0900 0000 0000 1147 1836, SWIFT: GIBASKBX

\n 12.	3.	Alžbeta	SMOLKOVÁ;
\n 17.	3.	Tamara	SMUTNÁ;
\n 20.	3.	Jana	MÁŤUŠOVÁ;
\n 25.	3.	Filip	KOVÁČ;
\n 26.	3.	Matej	KOLÁŘIK;
\n 27.	3.	Marta	MAJEROVÁ;
\n 28.	3.	Marta	BARKÓCI;
\n 29.	3.	Marcel	MAĎAR;

\narodeniny


\program{
\p  1 ; ne ;  9.30 ; Bohoslužby (J.~Szőllős + VP) ; 12.00 ; Spoločný obed v~SED ;
\p  2 ; po ; 18.00 ; Modlitebná skupinka ;.;;
\p  3 ; ut ; 15.00 ; Biblická hodina pre seniorov (P.~Pivka) ;.;;
\p  4 ; st ;.;;.;;
\p  5 ; št ; 18.00 ; Biblická hodina (F.~Barkóczi) ;.;;
\p  6 ; pi ; 17.30 ; Dorast ; 17.00 ; Svetový deň modlitieb (Cukrová~4) ;
\p  7 ; so ; 18.00 ; Mládež ;.;;
\p  8 ; ne ;  9.30 ; Bohoslužby (D.~M.~Chuchút) ;.;;
\p  9 ; po ; 18.00 ; Modlitebná skupinka / skupinka Rastieme v~Kristovi ;.;;
\p 10 ; ut ; 15.00 ; Biblická hodina pre seniorov (P.~Pivka) ;.;;
\p 11 ; st ;.;;.;;
\p 12 ; št ; 18.00 ; Pôstne zamyslenia 1. ;.;;
\p 13 ; pi ; 17.30 ; Dorast ;.;;
\p 14 ; so ; 18.00 ; Mládež ;.;;
\p 15 ; ne ;  9.30 ; Bohoslužby (R.~Nemec) ;.;;
\p 16 ; po ; 18.00 ; Modlitebná skupinka ;.;;
\p 17 ; ut ; 15.00 ; Biblická hodina pre seniorov (P.~Pivka) ;.;;
\p 18 ; st ;.;;.;;
\p 19 ; št ; 18.00 ; Pôstne zamyslenia 2. ;.;;
\p 20 ; pi ; 17.30 ; Dorast ;.;;
\p 21 ; so ; 18.00 ; Mládež ; 17.00 ; Veľkonočný koncert ;
\p 22 ; ne ;  9.30 ; Bohoslužby (F.~Barkóczi) ; 17.00 ; Veľkonočný koncert (Petržalka);
\p 23 ; po ; 18.00 ; Modlitebná skupinka / skupinka Rastieme v~Kristovi ;.;;
\p 24 ; ut ; 15.00 ; Biblická hodina pre seniorov (P.~Pivka) ;.;;
\p 25 ; st ; 17.30 ; Stretnutie sestier ;.;;
\p 26 ; št ; 18.00 ; Pôstne zamyslenia 3. ;.;;
\p 27 ; pi ; 17.30 ; Dorast ;.;;
\p 28 ; so ; 18.00 ; Mládež ;.;;
\p 29 ; ne ;  9.30 ; Bohoslužby (D.~M.~Chuchút) ;.;;
\p 30 ; po ; 18.00 ; Modlitebná skupinka ;.;;
\p 31 ; ut ; 15.00 ; Biblická hodina pre seniorov (P.~Pivka) ;.;;
}

\tiraz
\bye
